% !TeX root = ../概率论与数理统计笔记.tex
% 由主文档通过 \input 载入；请编译主文档。
\setcounter{chapter}{1} % 保留原书章号。
\chapter{条件概率、独立性与贝叶斯公式}
\label{chap:ps-conditioning}

本章研究获得新信息后概率如何变化，以及何时一个事件不会改变另一个事件的概率。
先修内容为第\ref{chap:ps-events}章；学习目标是正确选择条件、划分情形并辨认独立性假设。

\section{条件概率与乘法公式}
\label{sec:ps-conditional-probability}
定义\emph{条件概率}（conditional probability）$\mathbb P(A\mid B)$，首先限定 $\mathbb P(B)>0$。
推导乘法公式与多步试验的链式分解，说明树状分析的每条分支都带有相应条件。
零概率事件不能直接代入比值定义，其连续条件分布处理留待第\ref{chap:ps-multivariate}章。

\section{全概率公式与贝叶斯公式}
\label{sec:ps-total-bayes}
在有限或可数的事件划分下建立\emph{全概率公式}（law of total probability）与\emph{贝叶斯公式}（Bayes' formula），
明确互斥、覆盖与相关分母为正的条件。
由分层抽样和检测模型说明先验信息、条件概率与后验信息的区别，避免颠倒条件。

\section{事件的独立性}
\label{sec:ps-event-independence}
定义\emph{独立性}（independence）、两两独立和相互独立；对有限事件族要求相应各子族满足乘法关系。
证明独立性在取补运算下的保持性质，构造两两独立但不相互独立的例子。
比较互斥与独立，并说明非零概率的互斥事件不能独立。

\section{条件独立与条件化策略}
\label{sec:ps-conditional-independence}
定义给定正概率事件后的\emph{条件独立}（conditional independence），说明它与无条件独立之间一般不能相互推出。
介绍按第一步、来源类别或隐藏状态条件化的解题方法，并用分层比例说明总体关系与分组关系可能不同。

\section{例题、习题与学习检查}
规划双袋抽球、含假阳性与假阴性的检测、简单系统可靠性、两两独立反例等例题；检测例题使用明确给定的假想参数。
习题要求分别写出模型假设和推导步骤，并用直接计数核对小规模情形。
学习检查聚焦“条件是什么”“划分是否完备”“独立性从何而来”三个问题。






